\documentclass[UTF8,fontset=none,11pt,a4paper]{ctexart}
\usepackage{ipara-handbook}
\usepackage{amsmath,booktabs,tabularx,hyperref}
\usetikzlibrary{arrows.meta,positioning,fit,calc}

\handbooksetup{OS-B01 Wait / Block / Wakeup}{408 · OS Internal Bridge}{Condition · Queue · State · Reschedule}

\tikzset{
  qnode/.style={draw=iparaDeep,fill=iparaSoft,rounded corners=1.5mm,align=center,minimum width=2.6cm,minimum height=0.90cm,font=\small},
  evnode/.style={draw=iparaGreen,fill=white,rounded corners=1.5mm,align=center,minimum width=2.6cm,minimum height=0.90cm,font=\small},
  warnnode/.style={draw=iparaAccent,fill=white,rounded corners=1.5mm,align=center,minimum width=2.6cm,minimum height=0.90cm,font=\small},
  mainflow/.style={-{Latex[length=2mm]},thick,draw=iparaDeep},
  retflow/.style={-{Latex[length=2mm]},thick,dashed,draw=iparaGreen}
}

\begin{document}
\handbooktitle{Wait / Block / Wakeup}{条件不满足时，执行权怎样安全暂存、排队，并在事件后重新竞争 CPU}{408 · OS-B01}

\begin{corebox}{Mother Interface}
\[
\boxed{
\begin{aligned}
\text{Condition Unsatisfied} &\to \text{Register Wait + Block Safely} \\
&\to \text{Event / State Change} \to \text{Wake Eligible Waiter} \\
&\to \text{Ready/Runnable} \to \text{Scheduler/Dispatch} \\
&\to \text{Reacquire Protection + Recheck Condition}
\end{aligned}
}
\]
本 Bridge 专门拥有“同步条件、等待队列（Wait Queue）与任务状态（Task State）”之间的控制权交接契约；同步原语语义（OS-03）、设备中断处理（OS-05）和调度算法（OS-01/02）仍由各自专题独立拥有。
\end{corebox}

\begin{methodbox}{先定义接口两侧的词}
\begin{itemize}
  \item \kw{Condition Predicate（条件谓词）}：决定当前执行流能否继续推进的条件，如 \code{count > 0}、``request completed'' 或 ``lock available''。谓词由拥有共享状态的 Topic 定义，不能用“收到了 wakeup”替代。
  \item \kw{Wait Relation / Wait Queue（等待关系/等待集合）}：记录“哪些执行实体正在等待哪个条件/事件”。实现可以是链表、队列、哈希等待桶或其他结构；Bridge 只依赖\kw{等待者能够被事件定位}这一语义。
  \item \kw{Task State（任务执行状态）}：描述执行实体当前是否具备 CPU 调度资格。408 可使用 Ready/Running/Blocked；Linux 的 \code{TASK\_*} 状态只是一个工程实例，不能反过来当跨系统定义。
  \item \kw{Wakeup（唤醒）}：移除或更新等待关系，使满足资格的实体重新成为 Ready/Runnable。它改变的是\kw{调度资格}，不是“立即获得 CPU”，也不保证条件到真正运行时仍然成立。
\end{itemize}
\end{methodbox}

\begin{methodbox}{Owners 与职责边界}
\begin{itemize}
  \item \kw{左侧 Owner (OS-03 并发同步/OS-05 I/O)：}拥有锁、条件变量、信号量或外设 DMA 完成事件本体。
  \item \kw{右侧 Owner (OS-01/02 进程与调度)：}拥有 \code{task\_struct}、就绪队列（Run Queue）、CPU 时间片与上下文切换机制。
  \item \kw{本 Bridge Owns：}条件检查失败后“如何安全入队让出 CPU”、事件到达后“如何唤醒并重检条件”的协议链。
  \item \kw{Stops：}不重新定义死锁避免算法、PV 信号量数值含义或外设寄存器细节。
\end{itemize}
\end{methodbox}

\section{核心机制：为什么必须“先入队、改状态，后释放保护”}

在并发系统与操作系统内核中，条件检查与休眠之间存在极易触发的严重竞争态——\kw{丢失唤醒（Lost Wakeup）}。

若线程 A 先检查条件为假，在准备将自己标为 Blocked 之前发生上下文切换；此时线程 B 生产了资源并调用 \code{wakeup()}。由于线程 A 尚未挂入等待队列，线程 B 认为“当前没有等待者”因而直接返回。随后线程 A 恢复执行，盲目把自己标为 Blocked 并让出 CPU，陷入永久无法被唤醒的死锁状态！

\subsection{标准安全休眠与唤醒协议}
安全等待协议真正要求的不是固定 API 顺序，而是一个原子性不变量：\kw{从“条件仍为假”到“等待者已经可被未来事件找到并且不会继续消费 CPU”之间，不能出现 wakeup 无人接收的窗口。} 一种典型实现会：
\begin{enumerate}
  \item 在保护共享谓词的同步范围内检查条件；
  \item 建立当前执行实体与等待条件之间的可定位关系；
  \item 把执行实体变为不可运行/等待状态，并与释放保护、让出 CPU 做成不会丢失事件的协议；
  \item 进入调度路径，让其他 Ready/Runnable 实体继续推进。
\end{enumerate}
Linux wait queue、条件变量的 atomic release-and-wait、semaphore 的 P 操作都可以实现这份契约，但具体函数名、状态编码和内部队列结构不属于本 Bridge。

\section{母模型：阻塞排队与事件唤醒流转拓扑}

\begin{center}
\begin{tikzpicture}[node distance=0.75cm and 0.85cm]
  \node[qnode] (eval) {1. 检查条件\\$cond == false$};
  \node[qnode,right=of eval] (enq) {2. 登记等待\\建立 Wait Relation};
  \node[warnnode,right=of enq] (blk) {3. 变为不可运行\\安全让出 CPU};
  \node[evnode,below=1.1cm of blk] (event) {4. 状态/事件变化\\V / signal / I/O complete};
  \node[evnode,left=of event] (wake) {5. Wake Eligible\\恢复 Ready/Runnable};
  \node[qnode,left=of wake] (retry) {6. Scheduler/Dispatch\\重新获取保护并检查 $cond$};

  \draw[mainflow] (eval) -- (enq);
  \draw[mainflow] (enq) -- (blk);
  \draw[mainflow] (blk) -- (event);
  \draw[mainflow] (event) -- (wake);
  \draw[mainflow] (wake) -- (retry);
  \draw[retflow] (retry.north) -- node[left,font=\scriptsize]{若仍不满足再次入队} (eval.south);
\end{tikzpicture}
\end{center}

\subsection{为什么条件等待通常要在返回后重新检查谓词}
当线程从阻塞中被唤醒并真正重新获得 CPU 时，它面对的条件\textbf{不保证依然为真}：
\begin{itemize}
  \item \kw{多等待者竞争（典型 Mesa 风格语义）}：一个状态变化可能使若干等待者重新有资格竞争，但先运行者可能在后运行者取得锁之前再次改变共享状态；
  \item \kw{允许虚假/非唯一原因唤醒的接口}：某些条件等待 API 明确允许调用者在没有“一次通知对应一次永久真条件”的情况下返回。
\end{itemize}
因此，对\kw{条件变量/谓词式等待}，稳健模板是 \code{while(!condition) \{ wait(); \}}：wake 只意味着“值得再检查”。若题目给的是计数信号量、一次性完成事件等更强接口，应按该 Owner 的契约判断，不把 \code{while} 机械推广到所有等待原语。

\section{四种等待策略全景对比}

\begin{center}
\small
\renewcommand{\arraystretch}{1.25}
\begin{tabularx}{\linewidth}{>{\bfseries}p{0.18\linewidth} X X X}
\toprule
等待策略 & 核心机制与状态 & 优点与适用场景 & 劣势与工程代价 \\
\midrule
自旋等待 (Spin) & 执行实体仍可运行，在 CPU 上反复检查条件 & 避免睡眠/唤醒路径；当预计等待时间短于阻塞切换成本时可能划算 & 等待期间持续消耗 CPU；若持有者没有机会运行，继续自旋只会放大停滞 \\
阻塞休眠 (Blocking) & 建立等待关系并失去 CPU 调度资格，直到事件使其重新 Ready & CPU 可交给其他 runnable 工作；适合相对较长或由外部事件决定的等待 & 至少支付进入等待与以后重新运行的调度/现场管理成本，且有 wake-to-run 延迟 \\
让出执行 (Yield) & 当前实体主动放弃本次 CPU 机会，但通常没有声明自己在等待哪个确定事件 & 可降低原地独占，给其他 runnable 实体机会 & 仍可能很快再次被选中；不能替代真正的事件等待协议 \\
挂起 (Suspend) & 在教材模型中把“可运行/阻塞”再叠加一个暂不参与正常调度的维度 & 可表达外部暂停、负荷调节或经典换出场景 & Suspend 与 Blocked 正交；现代系统不应机械理解为“整个进程映像必然完整写入 swap” \\
\bottomrule
\end{tabularx}
\end{center}

\section{最小执行轨迹：生产者-消费者队列交接}

设有限缓冲区当前为空（\code{count == 0}）：
\begin{enumerate}
  \item \textbf{消费者 C1 到达}：获取互斥锁，检查 \code{count == 0} 为真；C1 将自身放入 \code{not\_empty} 等待队列，释放互斥锁并转入 \code{BLOCKED}；
  \item \textbf{生产者 P1 到达}：获取互斥锁，向缓冲区写入一个元素，\code{count = 1}；
  \item \textbf{P1 触发唤醒}：P1 在改变 \code{count} 后通知 \code{not\_empty} 等待者，C1 的等待关系被解除并重新成为 Ready/Runnable；
  \item \textbf{调度与重验}：P1 释放互斥锁并退出。调度器选中 C1，C1 恢复执行，首先\textbf{重新争抢互斥锁}；获取锁后重新检查 \code{count > 0} 成立，成功取出元素并减少 \code{count}。
\end{enumerate}

\begin{warnbox}{Anti-Bridge：三个必须阻断的错误直觉}
\begin{itemize}
  \item \(\text{Wakeup} \neq \text{Immediate Execution}\)：唤醒只是赋予调度资格，不代表立刻夺取正在运行任务的 CPU。
  \item \(\text{Device Interrupt} \neq \text{Task Rescheduled}\)：外设中断完成并唤醒等待线程后，CPU 仅将任务标记就绪，真正的上下文切换发生在该中断返回或下一个调度点。
  \item \(\text{Wakeup} \neq \text{Condition True}\)：被唤醒不代表资源已属于你，必须在重新获取锁后再次验证条件。
\end{itemize}
\end{warnbox}

\begin{corebox}{30 秒极速调用协议}
做题遇到“进程等待事件/信号量/I/O”时，严格按三步追踪：
\[
\boxed{\text{谁在等待哪个条件？}} \xrightarrow{\text{Event 使其 Ready/Runnable}} \boxed{\text{调度后按接口契约重新验证条件}}
\]
\end{corebox}

\end{document}
